%% schriftsatz.latex — Anpassbares Pandoc-Template für familienrechtliche Schriftsätze
%%
%% Pandoc-Variablen (via -V oder YAML-Frontmatter):
%%   mainfont      Schriftart             (Standard: Arial)
%%   fontsize      Schriftgröße           (Standard: 11pt)
%%   papersize     Papierformat           (Standard: a4)
%%   geometry      Seitenränder           (Standard: margin=2.5cm)
%%   linestretch   Zeilenabstand          (Standard: 1.2)
%%   lang          Sprache                (Standard: de-DE)
%%   briefkopf     Vorbereiteter LaTeX-Briefkopf (von generate-pdf.py erzeugt)
%%
%% ─────────────────────────────────────────────────────────────────────────────

\documentclass[
  11pt,
  a4paper,
]{article}

%% ── Sprache ──────────────────────────────────────────────────────────────────
\usepackage[ngerman]{babel}

%% ── Fonts (XeLaTeX) ──────────────────────────────────────────────────────────
\usepackage{fontspec}
\setmainfont{Arial}[
  BoldFont       = Arial Bold,
  ItalicFont     = Arial Italic,
  BoldItalicFont = Arial Bold Italic,
  Ligatures      = TeX,
]
\usepackage{microtype}

%% ── Seitenlayout ─────────────────────────────────────────────────────────────
\usepackage[left=3cm,right=2cm,top=2.5cm,bottom=2.5cm]{geometry}

%% ── Zeilenabstand & Absätze ──────────────────────────────────────────────────
\usepackage{setspace}
\setstretch{1.3}
\usepackage{parskip}
\setlength{\parskip}{6pt}
\setlength{\parindent}{0pt}

%% ── Überschriften ────────────────────────────────────────────────────────────
\usepackage{titlesec}
\titleformat{\section}{\normalfont\large\bfseries}{}{0em}{}
\titleformat{\subsection}{\normalfont\normalsize\bfseries}{}{0em}{}
\titleformat{\subsubsection}{\normalfont\normalsize\bfseries}{}{0em}{}
\titlespacing{\section}{0pt}{14pt}{6pt}
\titlespacing{\subsection}{0pt}{10pt}{4pt}
%% Keine automatische Abschnittsnummerierung
\setcounter{secnumdepth}{0}

%% ── Seitenzahlen ─────────────────────────────────────────────────────────────
\usepackage{fancyhdr}
\usepackage{lastpage}
\pagestyle{fancy}
\fancyhf{}
\fancyfoot[C]{\footnotesize Seite~\thepage~von~\pageref{LastPage}}
\renewcommand{\headrulewidth}{0pt}
\renewcommand{\footrulewidth}{0pt}

%% ── Tabellen ─────────────────────────────────────────────────────────────────
\usepackage{booktabs}
\usepackage{longtable}
\usepackage{array}
\usepackage{calc}
\usepackage{caption}
\captionsetup{font=small, labelfont=bf}
%% Pandoc 3.x setzt \LTcaptype auf "none" — LaTeX braucht diesen Counter:
\newcounter{none}

%% ── Listen ───────────────────────────────────────────────────────────────────
\usepackage{enumitem}
\setlist{noitemsep, topsep=4pt}

%% ── Sonstiges ────────────────────────────────────────────────────────────────
\usepackage{xcolor}
\usepackage[normalem]{ulem}
\usepackage[hidelinks]{hyperref}
\usepackage{needspace}

%% ── Pandoc-Kompatibilität ────────────────────────────────────────────────────
\providecommand{\tightlist}{%
  \setlength{\itemsep}{0pt}\setlength{\parskip}{0pt}}

%% ─────────────────────────────────────────────────────────────────────────────
\begin{document}

%% Briefkopf — von generate-pdf.py aus Markdown-Zeilen vor dem ersten --- erzeugt.
%% Zeilen mit "den \d" werden rechtsbündig gesetzt, Betreff fett+unterstrichen,
%% Aktenzeichen fett, alle anderen Zeilen linksbündig mit explizitem Zeilenumbruch.

\section{Notizen --- Az. 7 F 88/25}\label{notizen-az.-7-f-8825}

\begin{center}\rule{0.5\linewidth}{0.5pt}\end{center}

\subsection{Unverarbeitet}\label{unverarbeitet}

{\def\LTcaptype{none} % do not increment counter
\begin{longtable}[]{@{}
  >{\raggedright\arraybackslash}p{(\linewidth - 6\tabcolsep) * \real{0.2188}}
  >{\raggedright\arraybackslash}p{(\linewidth - 6\tabcolsep) * \real{0.2812}}
  >{\raggedright\arraybackslash}p{(\linewidth - 6\tabcolsep) * \real{0.2500}}
  >{\raggedright\arraybackslash}p{(\linewidth - 6\tabcolsep) * \real{0.2500}}@{}}
\toprule\noalign{}
\begin{minipage}[b]{\linewidth}\raggedright
Datum
\end{minipage} & \begin{minipage}[b]{\linewidth}\raggedright
Kontext
\end{minipage} & \begin{minipage}[b]{\linewidth}\raggedright
Inhalt
\end{minipage} & \begin{minipage}[b]{\linewidth}\raggedright
Status
\end{minipage} \\
\midrule\noalign{}
\endhead
\bottomrule\noalign{}
\endlastfoot
03.03.2026 & Telefonat Mandant & Finn hat letzte Woche in der KiTa
angeblich geweint beim Abholen durch AS --- Erzieherin hat nichts
Auffälliges berichtet, das steht in Widerspruch zum Antrag & offen \\
03.03.2026 & Telefonat Mandant & AS hat bei Übergabe am 28.02. erstmals
eigene Mutter mitgebracht --- Mandant fühlt sich überwacht, möchte das
nicht eskalieren & offen \\
\end{longtable}
}

\begin{center}\rule{0.5\linewidth}{0.5pt}\end{center}

\subsection{Verarbeitet}\label{verarbeitet}

{\def\LTcaptype{none} % do not increment counter
\begin{longtable}[]{@{}
  >{\raggedright\arraybackslash}p{(\linewidth - 6\tabcolsep) * \real{0.1186}}
  >{\raggedright\arraybackslash}p{(\linewidth - 6\tabcolsep) * \real{0.3729}}
  >{\raggedright\arraybackslash}p{(\linewidth - 6\tabcolsep) * \real{0.2712}}
  >{\raggedright\arraybackslash}p{(\linewidth - 6\tabcolsep) * \real{0.2373}}@{}}
\toprule\noalign{}
\begin{minipage}[b]{\linewidth}\raggedright
Datum
\end{minipage} & \begin{minipage}[b]{\linewidth}\raggedright
Inhalt (Kurzfassung)
\end{minipage} & \begin{minipage}[b]{\linewidth}\raggedright
Eingeordnet in
\end{minipage} & \begin{minipage}[b]{\linewidth}\raggedright
Entscheidung
\end{minipage} \\
\midrule\noalign{}
\endhead
\bottomrule\noalign{}
\endlastfoot
26.02.2026 & Mandant berichtet: Dienstreise Stuttgart war 16.10.2025,
nicht 15.10. wie AS behauptet & fakten.md / timeline.md & Widerspruch in
fakten.md eingetragen, Datumsfehler in Erwiderung VI. aufgegriffen \\
27.02.2026 & KiTa-Leiterin hat AG nach Entwicklungsgespräch spontan
gelobt: „Finn ist ein fröhliches Kind'' & nur-muendlich.md & Zitat zu
persönlich für Schriftsatz; als Redevorschlag in nur-muendlich.md \\
\end{longtable}
}

\end{document}
